\section{実験計画法1；実験計画法とは}
\label{lesson11}

\subsection{授業内容}
\subsubsection{科目の中でのこのコマの位置づけ}
\begin{tabular}[t]{ccccccc}
    記述統計[4] & 確率[2] & モデル[5] & 推測・推定[1] & 意思決定[0] & 実践[2]
\end{tabular}
\subsubsection{概要}
心理学の研究の多くは，従属変数に影響を及ぼす要因をコントロールし，統制群と比較することで平均因果効果を検証するところにある。心理学研究の具体例を基盤に，実験計画の種類，用語を確認し，平均差の効果は誤差に対する相対的な大きさで検討される原理を理解する。心理統計のテキストの多くは「分散分析」の文脈でこのことを論じるものがほとんどであるが，「線形モデルであること」「誤差と効果の対比であること」という二つの観点から読み解き，古典的テスト理論と同じ形の，関数で表現し誤差を分離するという心理統計の基本方針を見出さなければならない。
\subsubsection{コマ主題細目}
\begin{description}
\item[ランダム化]平均因果効果を検証するためには，ランダム化による比較，いわゆるランダム化比較実験(RCT,randomized controlled trial)が必要である。RCTによって誤差をコントロールすることができるようになることを理解する。
\begin{flushright}
    $\to$ \citet{西内201301}のP.101-107，\citet{大村201301}のP.17--40.
\end{flushright}

\item[「要因」と「水準」]実験計画法で用いられる用語，「要因」「水準」「群内計画」「群間計画」といった用語の使い方を整理する。特に心理学系の論文からどのような計画で行われた実験であったかを読み取ったり，自分自身が今後立案する場合はどのように考えれば良いか，具体的な例に即して理解する。
\begin{flushright}
    $\to$ \citet{Yamada2004}のP.174-177，\citet{大村201301}のP.48--52.
\end{flushright}

\item[回帰分析と実験計画]実験計画と回帰分析は深い関わりがある，というより数学的には同じ形式で表現されることを理解する。回帰分析と異なるのは，説明変数が連続的ではなく離散的である点であり，そのことが結果として散布図の形を変え，回帰線が平均差の比較になっていることを理解する。
\begin{flushright}
    $\to$ \citet{kosugi2019}のP.165--169
\end{flushright}

\item[平均差の効果]実験計画を一般線形モデルとして考えた場合，残差がどのように現れるか，また効果が大きい・小さいというのはどのように現れるかを理解する。平方和を分解($SS_t = SS_f+SS_e$)すること，誤差に対する比率で考えるという原理を理解する。
\begin{flushright}
    $\to$ \citet{大村201301}のP.84-88.
\end{flushright}


\end{description}
\subsubsection{キーワード}
\begin{itemize}
\item 実験計画
\item ランダム化比較実験
\item 平均因果効果
\item 一般線形モデル
\end{itemize}
\subsection{授業情報}
\paragraph{コマの展開方法}
講義
\paragraph{標準シラバスにおける位置づけ}
科目番号4；心理学研究法；2.データを用いた実証的な思考方法；C データの統計的記述\\
科目番号5；心理学統計法；1.心理学で用いられる統計手法；G 推測統計:代表値と散布度をめぐって(1)
\subsubsection{予習・復習課題}
\paragraph{予習}
また心理統計のテキストから，「分散分析」として説明されているもの(例えば\citet{Yamada2004}のP.162--165など)を一読しておくと良い。ただし，分散分析の文脈での議論は，最終的な効果の検証に当たって「帰無仮説検定」の枠組みを用いている。先んじてテキストを参照するものにとっては，$F分布$や自由度などの耳慣れない用語が出てくるが，この講義において帰無仮説検定について扱うのはまだ先になる。効果検証の判断基準については現段階で理解する必要はない。
\paragraph{復習}
既に習ったRの回帰分析関数をつかって，説明変数が離散的であるデータを使って分析を試み，どのような結果が出るかを確認すると良い。その際，説明変数がdata.frameの中でfactor型になっている必要が亜あることに注意する。水準は少ない方がわかりやすいため，2水準，3水準と徐々に増やして結果がどのように変化するかを確認しよう。分析の前にggplotで散布図を描くのを忘れないように。
